%
% tikztemplate.tex -- Template für standalone TIKZ Bilder
%
% (c) 2019 Prof Dr Andreas Müller, Hochschule Rapperswil
%
\documentclass[tikz,12pt]{standalone}
\usepackage{amsmath}
\usepackage{times}
\usepackage{txfonts}
\usepackage{pgfplots}
\usepackage{csvsimple}
\definecolor{darkgreen}{rgb}{0,0.6,0}
\definecolor{darkred}{rgb}{0.8,0,0}
\usetikzlibrary{arrows,intersections,math,calc}
\begin{document}
\def\skala{3}
\begin{tikzpicture}[>=latex,thick,scale=\skala]

\fill[color=gray!40] (0,0) rectangle (1,1);
\begin{scope}
\clip (0,-1) rectangle (3.1,2.2);
	\foreach \C in {0.05,0.1,0.2,...,2.2}{
		\draw[color=darkred,line width=1pt]
			plot[domain=0:3.1,samples=100] ({\x},{\C*exp(\x)});
		\draw[color=darkred,line width=1pt]
			plot[domain=0:3.1,samples=100] ({\x},{-\C*exp(\x)});
	}
	\foreach \D in {-0.8,-0.6,...,5.2}{
		\draw[color=darkgreen,line width=1pt]
			(0,\D) -- (3.2,{\D-3.2});
	}
	\draw[color=blue,line width=1pt] (0,2) -- (3,-1);
	\draw[color=blue,line width=1pt]
		plot[domain=0:1.8,samples=100] ({\x},{exp(\x-1)});
	\fill[color=blue] (1,1) circle[radius={0.07/\skala}];
\end{scope}
\draw[->] ({-0.1/\skala},0) -- (3.2,0) coordinate[label={$x$}];
\draw[->] (0,-1.1) -- (0,2.3) coordinate[label={left:$y$}];
\draw[color=darkred,line width=1pt] (0,0) -- (3.1,0);
\draw (0,-1) -- (3.2,-1);
\draw[color=blue,line width=2pt] (0,0) -- (0,2);
\draw[color=blue] ({-0.1/\skala},2) -- ({0.1/\skala},2);
\draw[color=blue] ({-0.1/\skala},0) -- ({0.1/\skala},0);
\node[color=blue] at (0,2) [left] {$2$};
\node at (0,1) [left] {$1$};
\node[color=blue] at (0,0) [left] {$0$};
\node at (1,0) [below] {$1$};
\draw (1,{-0.1/\skala}) -- (1,{0.1/\skala});


\end{tikzpicture}
\end{document}

